\documentclass[11pt,a4paper]{article}

\usepackage[utf8]{inputenc}
\usepackage[T1]{fontenc}
\usepackage[french]{babel}
\usepackage[top=2cm, bottom=2cm, left=2.2cm, right=2.2cm]{geometry}
\usepackage{setspace}
\setstretch{1.1}
\usepackage{lmodern}
\usepackage{microtype}
\usepackage[dvipsnames,table]{xcolor}
\usepackage{booktabs}
\usepackage{tabularx}
\usepackage{array}
\usepackage{listings}
\usepackage{amsmath, amssymb}
\usepackage[hidelinks]{hyperref}
\usepackage{fancyhdr}
\usepackage{lastpage}
\usepackage{tcolorbox}
\tcbuselibrary{skins, breakable, listings}
\usepackage{enumitem}
\usepackage{float}
\usepackage{multicol}
\usepackage{titlesec}
\usepackage{parskip}

% Couleurs
\definecolor{primary}{HTML}{1a5c38}
\definecolor{primarymid}{HTML}{27ae60}
\definecolor{primarylt}{HTML}{e8f5e9}
\definecolor{codegray}{HTML}{f4f4f4}
\definecolor{codecomment}{HTML}{7f8c8d}
\definecolor{codekw}{HTML}{8e44ad}
\definecolor{codestring}{HTML}{27ae60}
\definecolor{bordercolor}{HTML}{d0d7e3}
\definecolor{rowalt}{HTML}{f7f9fc}

% Titres compacts
\titlespacing*{\section}{0pt}{8pt}{4pt}
\titlespacing*{\subsection}{0pt}{5pt}{2pt}
\titleformat{\section}{\large\bfseries\color{primary}}{}{0em}{}[\color{primarymid}\titlerule]
\titleformat{\subsection}{\normalsize\bfseries\color{primary}}{}{0em}{}

% En-tête / pied
\pagestyle{fancy}
\fancyhf{}
\renewcommand{\headrulewidth}{0.3pt}
\renewcommand{\footrulewidth}{0.3pt}
\fancyhead[L]{\small\color{primary}\textbf{Axe 1 — Équations non-linéaires}}
\fancyhead[R]{\small\color{gray}Numerical Analysis App}
\fancyfoot[C]{\small\thepage\ / \pageref{LastPage}}

% Listings Python
\lstdefinestyle{py}{
  language=Python,
  basicstyle=\ttfamily\scriptsize,
  keywordstyle=\color{codekw}\bfseries,
  stringstyle=\color{codestring},
  commentstyle=\color{codecomment}\itshape,
  numbers=left, numberstyle=\tiny\color{gray}, numbersep=5pt,
  showstringspaces=false, breaklines=true, tabsize=4, frame=none,
  aboveskip=2pt, belowskip=2pt,
}

% Boîtes
\tcbset{
  code/.style={colback=codegray, colframe=bordercolor, arc=3pt,
               boxrule=0.5pt, left=4pt, right=4pt, top=2pt, bottom=2pt,
               fonttitle=\bfseries\scriptsize, title=#1, breakable},
  info/.style={colback=primarylt, colframe=primarymid, arc=3pt,
               boxrule=0.6pt, left=6pt, right=6pt, top=4pt, bottom=4pt},
}

\setlist[itemize]{noitemsep, topsep=2pt, leftmargin=*}
\setlist[enumerate]{noitemsep, topsep=2pt, leftmargin=*}

\begin{document}
\thispagestyle{fancy}

% ── Titre compact ─────────────────────────────────────────────────────────────
\begin{tcolorbox}[colback=primary, colframe=primary, arc=5pt, left=10pt, right=10pt,
                  top=6pt, bottom=6pt]
  \begin{minipage}{0.68\textwidth}
    {\large\bfseries\color{white} Rapport Technique — Axe 1}\\[2pt]
    {\small\color{white!80!black} Numerical Analysis App · Génie Logiciel 3\textsuperscript{ème} année · USTHB}
  \end{minipage}
  \hfill
  \begin{minipage}{0.28\textwidth}\raggedleft
    
    {\scriptsize\color{white!70!black} 2025--2028}
  \end{minipage}
\end{tcolorbox}

\vspace{4pt}

% ── 1. Introduction ───────────────────────────────────────────────────────────
\section{Introduction et périmètre}

L'Axe~1 couvre la \textbf{résolution numérique d'équations non-linéaires} : trouver
$x^*$ tel que $f(x^*)=0$, $f$ étant saisie par l'utilisateur. La contribution porte sur
quatre fichiers totalisant \textbf{1~282 lignes} :

\vspace{3pt}
\renewcommand{\arraystretch}{1.2}
\rowcolors{2}{rowalt}{white}
\begin{tabularx}{\linewidth}{>{\ttfamily\scriptsize}l X r}
\toprule
\normalfont\small\textbf{Fichier} & \small\textbf{Rôle} & \small\textbf{Lignes} \\
\midrule
axe1\_screen.py & Orchestration UI / algorithmes & 666 \\
ui/axe1\_widgets.py & Composants graphiques réutilisables & 392 \\
modules/chap\,1/nonlinear\_resolution.py & Algorithmes de résolution & 90 \\
modules/chap\,1/function\_analysis.py & Analyse mathématique de $f$ & 134 \\
\bottomrule
\end{tabularx}

% ── 2. Méthodologie ───────────────────────────────────────────────────────────
\section{Méthodologie de travail}

Le développement s'est organisé en quatre phases :

\begin{enumerate}
  \item \textbf{Audit de l'existant.} Le bouton \og{}Lancer\fg{} n'avait aucun
        \texttt{command=}, les quatre boutons d'analyse étaient inertes, la zone graphique
        n'affichait qu'un placeholder. Les algorithmes étaient implémentés mais jamais appelés.
  \item \textbf{Connexion UI~$\leftrightarrow$~Algorithmes.} Câblage de tous les boutons,
        parseur de formule utilisateur, gestion des erreurs à chaque étape.
  \item \textbf{Refactoring architectural.} Extraction des composants vers
        \texttt{ui/axe1\_widgets.py} et de la logique mathématique vers
        \texttt{function\_analysis.py} — principe \textit{Separation of Concerns}.
  \item \textbf{Nettoyage du module.} Suppression des fonctions \textit{hardcodées}
        (\texttt{f1\ldots f5}, \texttt{exo3\_function}) ; fonctions génériques paramétrées.
\end{enumerate}

Les modifications ont été versionnées sur GitHub dans la branche \texttt{work/axe1-screen},
fusionnée dans \texttt{main} via trois \textit{pull requests}.

% ── 3. Choix techniques ───────────────────────────────────────────────────────
\section{Choix techniques}

\subsection{Tkinter pur — zéro dépendance externe}

L'interface n'utilise ni \texttt{matplotlib} ni \texttt{ttk.Treeview} : le graphe est
tracé sur \texttt{tk.Canvas} par primitives vectorielles (\texttt{create\_line},
\texttt{create\_oval}) ; le tableau est dessiné cellule par cellule sur un
\texttt{Canvas} scrollable. Cela garantit la portabilité sur toute installation Python
standard sans étape d'installation.

\subsection{Parseur de formule sécurisé}

L'utilisateur saisit $f(x)$ en Python (ex. \texttt{x**3 - x - 2}). La fonction est
construite via \texttt{eval()} dans un environnement \textbf{isolé} : seules les
fonctions mathématiques explicitement autorisées sont exposées, \texttt{\_\_builtins\_\_}
est vidé.

\begin{tcolorbox}[code={Extrait — parseur sécurisé}]
\begin{lstlisting}[style=py]
SAFE_FUNCTIONS = {"np": np, "sin": np.sin, "exp": np.exp, ...}

def _parse_function(self):
    code = compile(self.func_entry.get(), "<f>", "eval")
    def f(x):
        env = dict(SAFE_FUNCTIONS); env["x"] = x
        return eval(code, {"__builtins__": {}}, env)
    return f
\end{lstlisting}
\end{tcolorbox}

\subsection{Dérivée numérique automatique}

Pour Newton, $f'$ est calculée automatiquement par différences finies centrées :
$f'(x)\approx\dfrac{f(x+h)-f(x-h)}{2h}$, $h=10^{-7}$ — sans saisie supplémentaire.

% ── 4. Algorithmes ────────────────────────────────────────────────────────────
\section{Algorithmes implémentés}

\begin{multicols}{2}

\subsection{Dichotomie (Bissection)}

Si $f$ est continue sur $[a,b]$ et $f(a)\cdot f(b)<0$, on pose :
\[
  m_n = \frac{a_n+b_n}{2}
\]
puis on restreint l'intervalle selon le signe de $f(a_n)\cdot f(m_n)$.
Convergence \textbf{garantie}, erreur $\leq (b-a)/2^n$.

\columnbreak

\subsection{Newton--Raphson}

\[
  x_{n+1} = x_n - \frac{f(x_n)}{f'(x_n)}
\]
Convergence \textbf{quadratique} au voisinage d'une racine simple.
3~à~5~itérations suffisent pour $\varepsilon=10^{-8}$.

\end{multicols}

\subsection{Point fixe}

On itère $x_{n+1}=g(x_n)$ avec $g(x)=x-f(x)$. Converge si $g$ est
\textbf{stable} ($g([a,b])\subseteq[a,b]$) et \textbf{contractante}
($k=\max|g'(x)|<1$). L'application vérifie ces deux conditions à la demande.

% ── 5. Tests ──────────────────────────────────────────────────────────────────
\section{Tests et validations}

\subsection{Cas nominaux}

\renewcommand{\arraystretch}{1.25}
\rowcolors{2}{rowalt}{white}
\begin{tabularx}{\linewidth}{l >{\ttfamily}l l c c}
\toprule
\textbf{Algo} & \normalfont\textbf{f(x)} & \textbf{[a,b] / x$_0$}
  & \textbf{Racine} & \textbf{Iter.} \\
\midrule
Dichotomie  & x**3 - x - 2    & $[1,2]$    & 1.52137971 & 20 \\
Dichotomie  & exp(x) + x - 2  & $[0,1]$    & 0.44285440 & 20 \\
Newton      & x**3 - x - 2    & $x_0=1.5$  & 1.52137971 &  4 \\
Newton      & exp(x) + x - 2  & $x_0=0.5$  & 0.44285440 &  3 \\
Point Fixe  & x - 2**(-x)     & $x_0=0.5$  & 0.64118539 & 17 \\
\bottomrule
\end{tabularx}

\vspace{4pt}

\subsection{Cas limites}

\begin{tcolorbox}[info]
\renewcommand{\arraystretch}{1.15}
\begin{tabularx}{\linewidth}{l X}
\textbf{Cas} & \textbf{Comportement vérifié} \\
\hline
$f(a)\cdot f(b)\geq 0$ & Message d'erreur contextuel, pas de crash \\
$f'(x_0)=0$ (Newton)   & Arrêt + dialog \og{}dérivée nulle\fg{} \\
Expression invalide     & Signalée avant toute exécution \\
$a\geq b$               & Erreur \og{}Il faut $a<b$\fg{} \\
Point fixe divergent    & Badge \og{}non convergé\fg{} + avertissement \\
\end{tabularx}
\end{tcolorbox}

% ── 6. Conclusion ─────────────────────────────────────────────────────────────
\section{Conclusion}

L'Axe~1 livre un module complet et autonome : interface \textbf{entièrement fonctionnelle},
\textbf{zéro dépendance externe}, architecture en couches stricte, saisie dynamique de
$f(x)$, et robustesse sur tous les cas d'erreur. Les trois algorithmes (Dichotomie, Newton,
Point Fixe) sont validés sur des fonctions à racines connues analytiquement.

\end{document}